% 第5章图：梯子形图（上排1,2,3，下排4,5,6，横档14,25,36）的树分解。
% 袋 {1,2,4}-{2,4,5}-{2,3,5}-{3,5,6} 自左向右扫过梯子；含2的袋（蓝色）连成一段，相邻袋的交集是扫描时要记住的接口。
\begin{tikzpicture}[x=1mm,y=1mm,>=Stealth,
  v/.style={circle,draw=FigureStroke,fill=white,line width=.7pt,minimum size=5.4mm,inner sep=0pt},
  vt/.style={v,draw=FigureBlue,fill=FigureBlueFill,line width=1.1pt},
  e/.style={FigureStroke,line width=.7pt},
  bag/.style={rectangle,rounded corners=1.8pt,draw=FigureStroke,fill=white,line width=.7pt,minimum width=12mm,minimum height=6.5mm,inner sep=1.5pt},
  has2/.style={bag,draw=FigureBlue,fill=FigureBlueFill,line width=1.1pt},
  sc/.style={font=\figurefont\footnotesize,text=FigureInk},
  itf/.style={font=\figurefont\footnotesize,text=FigureGreen},
  note/.style={font=\FigureNoteFont\fontsize{7.5}{9}\selectfont,text=FigureMuted}]
% (a) ladder graph
\begin{scope}
 \node[v] (g1) at (0,14) {$1$}; \node[vt] (g2) at (10,14) {$2$}; \node[v] (g3) at (20,14) {$3$};
 \node[v] (g4) at (0,0) {$4$}; \node[v] (g5) at (10,0) {$5$}; \node[v] (g6) at (20,0) {$6$};
 \draw[e] (g1)--(g2)--(g3) (g4)--(g5)--(g6) (g1)--(g4) (g2)--(g5) (g3)--(g6);
 \node[sc] at (10,-10) {(a) 梯子形图};
\end{scope}
% (b) path decomposition
\begin{scope}[xshift=33mm,yshift=7mm]
 \node[has2] (b1) at (0,0) {$\{1,2,4\}$};
 \node[has2] (b2) at (19.5,0) {$\{2,4,5\}$};
 \node[has2] (b3) at (39,0) {$\{2,3,5\}$};
 \node[bag]  (b4) at (58.5,0) {$\{3,5,6\}$};
 \draw[FigureStroke,line width=.8pt] (b1)--(b2) node[midway,above=1.2mm,itf]{$\{2,4\}$};
 \draw[FigureStroke,line width=.8pt] (b2)--(b3) node[midway,above=1.2mm,itf]{$\{2,5\}$};
 \draw[FigureStroke,line width=.8pt] (b3)--(b4) node[midway,above=1.2mm,itf]{$\{3,5\}$};
 \draw[FigureBlue,line width=.8pt] (-6,-6) -- (-6,-7.5) -- (45,-7.5) -- (45,-6);
 \node[note,text=FigureBlue] at (19.5,-10.5) {含$2$的袋：$T_2$连成一段};
 \node[sc] at (29.25,-17) {(b) 宽度为$2$的树分解};
\end{scope}
\end{tikzpicture}
